\section{Objectifs}

\begin{itemize}
\item Calculer avec des puissances entières relatives, des racines carrées et des expressions fractionnaires simples.
\item Développer, réduire et factoriser des expressions simples.
\item Utiliser les identités remarquables et choisir la forme d'une expression adaptée au problème.
\item Comparer deux quantités par différence ou, lorsqu'elles sont strictement positives, par quotient.
\item Manipuler des inégalités en tenant compte du signe du nombre par lequel on multiplie.
\item Exprimer une variable en fonction des autres dans une relation simple.
\end{itemize}

\section{Prérequis}

\begin{itemize}
\item priorités opératoires
\item développement simple
\item fractions et puissances
\end{itemize}

\section{Activité de découverte}

Pour comparer deux offres de location, on écrit le prix total avec une variable \texttt{x}. Selon la question, une forme développée, réduite ou factorisée peut être plus efficace. Le calcul littéral sert ici à choisir une stratégie, pas seulement à transformer des écritures.

\section{Vocabulaire}

\begin{itemize}
\item \textbf{expression}, \textbf{développer}, \textbf{factoriser}, \textbf{réduire}
\item \textbf{puissance entière relative}, \textbf{racine carrée}, \textbf{expression fractionnaire}
\item \textbf{forme développée}, \textbf{forme factorisée}, \textbf{forme réduite}
\item \textbf{comparaison additive}, \textbf{comparaison multiplicative}
\end{itemize}

\section{Cours}

\subsection{Puissances et racines carrées}

On réinvestit les règles de calcul sur les puissances entières relatives et sur les racines carrées. Pour tout réel \texttt{a}, \texttt{√(a²)=|a|}. Pour \texttt{a ≥ 0} et \texttt{b ≥ 0}, \texttt{√(ab)=√a × √b}.

Une expression fractionnaire n'est définie que lorsque son dénominateur est non nul. Cette condition doit être déterminée avant tout calcul ou simplification.

\subsection{Calcul littéral}

Les identités remarquables à maîtriser sont :

\begin{itemize}
\item \texttt{(a+b)² = a² + 2ab + b²} ;
\item \texttt{(a-b)² = a² - 2ab + b²} ;
\item \texttt{(a+b)(a-b) = a²-b²}.
\end{itemize}

On factorise aussi des expressions simples en mettant un facteur commun en évidence, par exemple \texttt{ax²+bx = x(ax+b)} et \texttt{ax+bx=(a+b)x}.

Le choix de la forme dépend de la tâche :

\begin{itemize}
\item développer et réduire pour effectuer certains calculs ou comparer ;
\item factoriser pour résoudre une équation produit ou étudier un signe ;
\item conserver une forme fractionnaire lorsqu'elle rend une structure visible.
\end{itemize}

\subsection{Inégalités et comparaisons}

\begin{itemize}
\item On peut additionner deux inégalités de même sens lorsque les termes sont comparables.
\item Multiplier une inégalité par un réel positif conserve son sens.
\item Multiplier une inégalité par un réel négatif inverse son sens.
\item Deux quantités peuvent être comparées par leur \textbf{différence} ; si elles sont strictement positives, on peut aussi comparer leur \textbf{quotient}.
\end{itemize}

Selon le contexte, ces deux comparaisons correspondent à une variation additive ou multiplicative.

\subsection{Isoler une variable}

Dans une relation comportant plusieurs variables, on peut exprimer l'une d'elles en fonction des autres. Exemples :

\begin{itemize}
\item \texttt{U = RI} donne \texttt{I = U/R} si \texttt{R ≠ 0} ;
\item \texttt{d = vt} donne \texttt{v = d/t} si \texttt{t ≠ 0} ;
\item \texttt{ax + by = c} donne \texttt{y = (c-ax)/b} si \texttt{b ≠ 0}.
\end{itemize}

\section{Exemples corrigés}

\begin{itemize}
\item \texttt{(x+3)² = x²+6x+9}.
\item \texttt{5x²-10x = 5x(x-2)}.
\item \texttt{√72 = √(36×2)=6√2}.
\item Pour \texttt{a=3} et \texttt{b=5}, \texttt{a-b=-2<0}, donc \texttt{a<b}.
\item Si \texttt{x<4}, alors \texttt{-2x>-8} car la multiplication par \texttt{-2} inverse le sens.
\end{itemize}

\coursefigure{/images/mathematiques/latex-migration/2nde-calcul-litteral-puissances-racines-figure-1.svg}{Trois formes d'une expression reliées à trois usages.}{La forme d'une expression est choisie en fonction de l'objectif.}

\section{Algorithmique en Python}

\subsection{Première puissance dépassant un seuil}

\begin{verbatim}
def premiere_puissance(base, seuil):
    n = 0
    valeur = 1
    while valeur < seuil:
        n += 1
        valeur *= base
    return n, valeur

print(premiere_puissance(2, 1000))
\end{verbatim}

La boucle détermine la première puissance du nombre positif choisi qui atteint ou dépasse le seuil.

\section{Erreurs fréquentes}

\begin{itemize}
\item Écrire à tort \texttt{(a+b)²=a²+b²}.
\item Oublier que \texttt{√(a²)=|a|}, et non toujours \texttt{a}.
\item Simplifier une expression fractionnaire sans exclure les valeurs qui annulent le dénominateur.
\item Oublier d'inverser le sens d'une inégalité après multiplication ou division par un nombre négatif.
\item Transformer une expression sans tenir compte de l'objectif du problème.
\end{itemize}

\section{Formules et idées à retenir}

\begin{itemize}
\item \texttt{(a+b)²=a²+2ab+b²}
\item \texttt{(a-b)²=a²-2ab+b²}
\item \texttt{(a+b)(a-b)=a²-b²}
\item \texttt{√(a²)=|a|}
\item \texttt{√(ab)=√a√b} pour \texttt{a,b≥0}
\end{itemize}

\section{Synthèse}

Le calcul algébrique de seconde ne consiste pas à appliquer une liste de recettes : il faut savoir choisir une écriture, gérer les conditions de définition, comparer des quantités et isoler une variable pour résoudre ou modéliser un problème.
